% Problem-section file following ../_template.tex.
\section{Resources for implementing Gibbs-preserving channels}
\label{sec:problem-37}

\paragraph{Problem.}
What are the exact one-shot coherence and work costs of implementing an
arbitrary Gibbs-preserving channel by thermal operations?  Let finite systems
$S$ and $S'$ have Hamiltonians $H_S$ and $H_{S'}$ at inverse temperature
$\beta$, with Gibbs states
\begin{equation}
  \gamma_X:=\frac{e^{-\beta H_X}}{\operatorname{Tr}(e^{-\beta H_X})},
  \qquad X\in\{S,S'\}.
  \label{eq:p37-gibbs-states}
\end{equation}
A channel $\Phi:S\to S'$ is Gibbs preserving when it maps the first state in
Eq.~\eqref{eq:p37-gibbs-states} to the second.  For approximation error
$\varepsilon\geq0$, define its quantum-Fisher-information coherence cost by
\begin{equation}
  F_c^{\varepsilon}(\Phi)
  :=\inf_{\substack{(R,H_R),\ \eta_R\in\mathcal D(R)\\
                    \mathcal T:S\otimes R\to S'\ \mathrm{thermal}}}
  \left\{
    F_Q(\eta_R,H_R):
    D_{\mathrm{ch}}\!\left(
      \Phi,\mathcal T(\,\cdot\,\otimes\eta_R)
    \right)\leq\varepsilon
  \right\},
  \label{eq:p37-coherence-cost}
\end{equation}
where the infimum ranges over finite resource systems $R$ with Hamiltonian
$H_R$, and $\mathcal T$ is a thermal operation implemented with Gibbs
ancillas and an energy-conserving unitary; set $\inf\varnothing:=+\infty$.
The distance $D_{\mathrm{ch}}$ is the supremum
of purified distance between the two channel outputs over inputs entangled
with an arbitrary reference.  If
$\eta_R=\sum_j\lambda_j\lvert j\rangle\!\langle j\rvert$, the quantity in
Eq.~\eqref{eq:p37-coherence-cost} is
\begin{equation}
  F_Q(\eta_R,H_R)
  :=2\!\sum_{j,k:\,\lambda_j+\lambda_k>0}
  \frac{(\lambda_j-\lambda_k)^2}{\lambda_j+\lambda_k}
  \left\lvert\langle j\rvert H_R\lvert k\rangle\right\rvert^2.
  \label{eq:p37-quantum-fisher-information}
\end{equation}
Equation~\eqref{eq:p37-quantum-fisher-information} fixes the normalization of
the coherence measure used in Eq.~\eqref{eq:p37-coherence-cost}.
Determine Eq.~\eqref{eq:p37-coherence-cost}, up to matching bounds, for every
Gibbs-preserving $\Phi$ and every $\varepsilon$.  Give the analogous tight
deterministic work cost when an ideal battery begins and ends in sharp energy
states and the battery's energy loss is charged as work.  Which structural
features of $\Phi$ determine whether the exact costs are finite, and what is
their sharp divergence as $\varepsilon\downarrow0$ when they are infinite at
zero error?

\subsection*{Status}

Unsolved

\subsection*{Source}

Tajima and Takagi explicitly leave a tight channel-by-channel
characterization of coherence and work costs for Gibbs-preserving operations
open \sourcecite{ref:p37-tajima}{TT25}.

\subsection*{Progress}

\begin{itemize}
  \item Every thermal operation is Gibbs preserving and time-translation
  covariant, whereas a general Gibbs-preserving channel need not be covariant
  and may create energy coherence.  The operation classes are therefore
  inequivalent \sourcecite{ref:p37-faist-oppenheim}{FOR15}.

  \item Universal optimal work rates are known in the independent and
  identically distributed limit, and one-shot implementations are known for
  important time-covariant channels
  \sourcecite{ref:p37-faist-berta}{FBB21}.  These results do not determine
  either cost requested above for every Gibbs-preserving channel.

  \item Some Gibbs-preserving channels cannot be implemented exactly by a
  thermal operation aided by any finite amount of coherence.  For pairwise
  reversible channels, lower bounds on
  Eq.~\eqref{eq:p37-coherence-cost} scale as $1/\varepsilon^2$ and are
  essentially attained by suitable examples
  \sourcecite{ref:p37-tajima}{TT25}.  This disproves universal finite
  coherence sufficiency but does not give a tight formula for arbitrary
  channels.
\end{itemize}

\subsection*{References}

\begin{enumerate}[label={},leftmargin=4.8em,labelsep=0.5em]
  \footnotesize
  \item[\textup{[FOR15]}]\label{ref:p37-faist-oppenheim}
  P. Faist, J. Oppenheim, and R. Renner, ``Gibbs-Preserving Maps Outperform
  Thermal Operations in the Quantum Regime,''
  \emph{New Journal of Physics} \textbf{17}, 043003 (2015).
  \href{https://doi.org/10.1088/1367-2630/17/4/043003}%
       {doi:10.1088/1367-2630/17/4/043003};
  \href{https://arxiv.org/abs/1406.3618}{arXiv:1406.3618}.

  \item[\textup{[FBB21]}]\label{ref:p37-faist-berta}
  P. Faist, M. Berta, and F. G. S. L. Brand\~ao,
  ``Thermodynamic Implementations of Quantum Processes,''
  \emph{Communications in Mathematical Physics} \textbf{384}, 1709--1750
  (2021). \href{https://doi.org/10.1007/s00220-021-04107-w}%
       {doi:10.1007/s00220-021-04107-w};
  \href{https://arxiv.org/abs/1911.05563}{arXiv:1911.05563}.

  \item[\textup{[TT25]}]\label{ref:p37-tajima}
  H. Tajima and R. Takagi, ``Gibbs-Preserving Operations Requiring Infinite
  Amount of Quantum Coherence,'' \emph{Physical Review Letters} \textbf{134},
  170201 (2025).
  \href{https://doi.org/10.1103/PhysRevLett.134.170201}%
       {doi:10.1103/PhysRevLett.134.170201};
  \href{https://arxiv.org/abs/2404.03479}{arXiv:2404.03479}.
\end{enumerate}

\subsection*{Comment}

The conclusion of \sourcecite{ref:p37-tajima}{TT25} explicitly leaves a
complete tight characterization of coherence cost for all Gibbs-preserving
operations and the corresponding work-cost problem.  Universal exact
implementation with finite coherence is already known to be false; the open
task is the channel-by-channel cost characterization.

\subsection*{Tag}

Quantum thermodynamics; Thermal operations; Quantum resource theories;
Asymmetry and quantum reference frames; Quantum channels;
One-shot quantum information

\subsection*{ID}

\texttt{op\_56578531a2c610fe}
